% --- Άσκηση 13171 (13171.tex) ---
\Askhsh{13171}{4}
\begin{ylh}{2}
    \item 2.1. Οι Πράξεις και οι Ιδιότητές τους
    \item 5.1. Ακολουθίες
    \item 5.2. Αριθμητική πρόοδος
\end{ylh}
Το άθροισμα των $\nu$ πρώτων διαδοχικών όρων μιας ακολουθίας $(a_{\nu})$ είναι $a_{1} + a_{2} + \ldots + a_{\nu} = S_{v}\  = \ 2v^{2}\  + \ 3\nu$ , $v \in N$ με $\nu \geq 1$.
\begin{alist}
\item Να βρείτε τον πρώτο όρο $a_{1}$.
\monades{5}
\item Να αποδείξετε ότι $S_{v - 1}\  = \ 2v^{2} - \nu - 1$, $v \geq 2$
\monades{6}
\item Να αποδείξετε ότι $a_{v} = 4v + 1$, $v \geq 1$
\monades{7}
\item Να αποδείξετε ότι αυτή η ακολουθία είναι αριθμητική πρόοδος, της οποίας να βρείτε τη διαφορά $\omega$.

\monades{7}
\end{alist}
